To examine whether the choice depends on the selection priorities, an additional analysis uses the same eight-candidate validation data, varying the minimum per-run pass fraction over 80\%, 90\%, and 100\%, and comparing dimension-first with energy-first ranking. Changing the pass-fraction requirement does not change any choice at the tested error tolerances. The ranking preference matters when only Node 2 is allowed at $\varepsilon=0.10$: dimension-first selects $\mathcal O_{23}$, whereas energy-first selects full observations. All other choices are unchanged. This analysis uses validation data only and does not alter the configurations already chosen for independent testing. The subsequently completed paired subsets are not retroactively inserted into the original candidate set or the previously selected configurations.

We then examine whether the nominal choices retain their requested accuracy when response parameters change. The original selected Node 2/$\mathcal O_{23}$ and Node 3/$\mathcal O_{13}$ policies are evaluated without changing their parameters at their respective tolerances 0.10 and 0.05. For 40 new common initial conditions, response weights are set to $W_r=W\odot(1+\eta U)$, where each entry of $U$ is sampled uniformly from $[-1,1]$ and $\eta\in\{0,0.01,0.03,0.05\}$. The drive system is unchanged. Each initial condition uses the same mismatch direction across amplitudes and policies; an independent random stream preserves identical initial states. Zero weights remain zero. No retraining or configuration reselection is performed.

\begin{table*}[pos=!t,width=\textwidth]
\centering
\caption{Performance of the selected configurations under response-weight mismatch, without retraining or reselection. Pass fractions are the minimum across three training runs, each evaluated on the same 40 new initial conditions. Mean tail error $z$ and energy $E_u$ average runs and initial conditions. Pass fractions track each configuration against its own previously fixed tolerance.}
\label{tab:configuration_transfer}
\small
\begin{tabular}{cccccc}
\toprule
Configuration & $\varepsilon$ & $\eta$ & Minimum pass & Mean $z$ & Mean $E_u$\\
\midrule
$2/\mathcal O_{23}$ & 0.10 & 0\% & 100.0\% & 0.0266 & 10.08\\
$2/\mathcal O_{23}$ & 0.10 & 1\% & 97.5\% & 0.0402 & 10.11\\
$2/\mathcal O_{23}$ & 0.10 & 3\% & 57.5\% & 0.0942 & 10.20\\
$2/\mathcal O_{23}$ & 0.10 & 5\% & 35.0\% & 0.1571 & 10.39\\
$3/\mathcal O_{13}$ & 0.05 & 0\% & 100.0\% & 0.0084 & 5.84\\
$3/\mathcal O_{13}$ & 0.05 & 1\% & 97.5\% & 0.0181 & 5.85\\
$3/\mathcal O_{13}$ & 0.05 & 3\% & 75.0\% & 0.0455 & 5.91\\
$3/\mathcal O_{13}$ & 0.05 & 5\% & 42.5\% & 0.0743 & 6.01\\
\bottomrule
\end{tabular}
\end{table*}

Table~\ref{tab:configuration_transfer} shows that both configurations retain a minimum pass fraction of 97.5\% at 1\% relative mismatch, but fall below the original 90\% requirement at 3\% and 5\%. No state-protection boundary is reached in these 960 trajectories. The drop in pass fraction identifies the sensitivity of the selected configurations to larger response-weight deviations and motivates incorporating the intended parameter range into configuration validation.
